\section{Kết luận}

Dãy ước lượng tách chặn hội tụ thành tốc độ nền $\beta_k$ và thành phần sai số tích lũy $\delta_k$. Với IAPPA1, sai số được khuếch đại bởi $1/\alpha_k$, cộng dồn qua $\eta_k$, rồi các tổng riêng được bình phương. Với IAPPA2, quan hệ dưới vi phân gần đúng được xác định trực tiếp và $\delta_k$ chỉ chứa tổng có trọng số của các bình phương sai số.

IAPPA1 bảo đảm hội tụ khi $q>3/2$, nhưng chặn tổng quát không bảo toàn tốc độ $O(k^{-2})$ khi có sai số loại 1. IAPPA2 bảo đảm hội tụ khi $q>1/2$ và giữ tốc độ $O(k^{-2})$ khi $q>3/2$. Vì vậy, chất lượng của phép tính gần đúng phải được đánh giá đồng thời qua độ lớn sai số và cấu trúc của chứng chỉ dừng.

